\noindent\textbf{Scope of improvement.}\quad
\toolname targets patch-related regression checking, not whole-project bug finding or vulnerability-class discovery.  A checker may know \emph{where} to warn yet miss \emph{why} the warning should disappear after the fix.  The refinement target is a change that reduces fixed-side noise while keeping the vulnerable trigger intact on the paired patch-related objects.  Use on a changed fork or later version requires additional stability tests; paired PDS alone does not justify that deployment.

\noindent\textbf{Why post-synthesis and LLM-generated checkers?}\quad
The patch, paired revisions, and an existing checker give a fixed artifact and an observable fixed-noisy failure to repair.  Analyzer evidence might also inform a fresh generator or a handwritten checker, but those settings change the starting artifact or authoring workflow and have not been evaluated here.  We separate generation from refinement to attribute any patch-local false-positive reduction to an edit of the same checker, not to a new generation draw; the method is not inherently limited to LLM authorship if an existing checker and paired validation target are supplied.

\noindent\textbf{Evidence as context compression.}\quad
The relevant semantic condition may span multiple files and analyzer outputs.  Passing the full repository to the model is infeasible, while passing only the patch is too sparse.  \toolname bridges this gap by collecting a provenance-tagged summary of patch-relevant facts, compact enough for the agent yet rich enough to guide edits.  This positions checker refinement closer to specification completion than to report suppression~\cite{clarke2000cegar,solarlezama2006sketch,nguyen2013semfix,mechtaev2016angelix,long2016prophet}.

\noindent\textbf{Backend complementarity.}\quad
E2 exercises both CSA and CodeQL detector formats, but does not by itself show that every evidence role was populated from an analyzer-internal representation.  In the audited CSA cohort, the verified native interface is a CFG/call-graph dump; source-derived and diagnostic records remain separately labelled.  A CodeQL-native coverage claim requires retained database-query provenance.

\noindent\textbf{Failure modes.}\quad
The dominant failure mode is over-pruning: a candidate suppresses fixed-side warnings by deleting or over-constraining the vulnerable trigger.  PDS rejects such cases as validity loss, never success.  Other modes include evidence insufficiency (a collector misses the needed fact), backend mismatch (the analyzer cannot expose the required relation), and model-side edit quality (the model cannot turn a correct fact into compilable code).  In each case the gate rejects the candidate rather than silently accepting it.

\noindent\textbf{Limited semantic generalization.}\quad
\toolname validates patch-local discrimination, not all instances of a vulnerability class~\cite{chakraborty2021deep,pearce2022copilot,perry2023users,khoury2023chatgpt}.  Some edits introduce a structural predicate; others, including the motivating AST-body edit, are tied to a patch function.  A post-hoc case-19 probe survives insertion of an unrelated statement and renaming the enclosing function, but misses a consistently renamed cleanup callback.  Thus even an automatic PDS candidate may remain callback-specific; this diagnostic is not a held-out transfer result.

\noindent\textbf{Precision and scaling boundary.}\quad
The strict gate protects vulnerable-side recall while reducing fixed-side alerts within the patch-related validation scope.  Across three matched repeats on the 12 selected cases, native \toolname retains 23, 23, and 18 of 37 starting fixed alerts, below KNighter's paired 28, 27, and 22; all accepted native candidates remain vulnerable-side hits.  The frozen CSA evidence replay completed on all 12 cases with four to nine records per bundle (70 total), and portfolio arithmetic retains all 39 checkers under PDS-only adoption without sending whole repositories to the editor.  This is bounded, automated workflow scale in the measured cohort, not whole-kernel throughput or broad vulnerability-class generalization.
